\chapter{Conclusiones}

\begin{enumerate}
  \item Sin protección, la corrida crítica alcanzó 619 procesos, 100\% de CPU
        y 73\% de RAM.
  \item Anti-Bunny Virus detectó siete PID y escaló sus alertas de sospechosas
        a críticas tras dos ciclos, con un máximo de 30,78 hijos/s.
  \item Las pruebas automatizadas confirmaron el rechazo de objetivos no
        autorizados y de identidades inválidas.
  \item La hipótesis quedó parcialmente respaldada: se validó la detección,
        pero no la contención experimental, la recuperación ni la sobrecarga.
\end{enumerate}

\section{Cumplimiento de objetivos}

\begin{table}[H]
  \centering
  \caption{Cumplimiento de los objetivos del proyecto.}
  \label{tab:cumplimiento-objetivos}
  \small
  \begin{tabular}{p{0.13\textwidth}p{0.18\textwidth}p{0.57\textwidth}}
    \toprule
    Objetivo & Estado & Evidencia \\
    \midrule
    General & Parcial & Prototipo implementado y detección de procesos demostrada; faltan pruebas experimentales de memoria, archivos y contención. \\
    OE1 & Parcial & El sistema recolecta PID, PPID, PGID, CPU, RSS y archivos; la evidencia entregada solo valida procesos. \\
    OE2 & Parcial & Siete PID escalaron de sospechosos a críticos tras dos ciclos, con máximo de 30,78 hijos/s. \\
    OE3 & Parcial & Las pruebas automatizadas validaron los rechazos de seguridad; no se registró terminación real de un PGID autorizado. \\
    OE4 & Parcial & Se documentaron dos corridas sin monitor y una con detección; faltan condiciones equivalentes, tres repeticiones, contención y sobrecarga. \\
    \bottomrule
  \end{tabular}
\end{table}
